\chapter{Estrutura Geral do TCC}

O TCC deve ser organizado em elementos pré-textuais, textuais e pós-textuais. Essa organização facilita a identificação do trabalho, a leitura do conteúdo, a localização de partes específicas e a conferência das fontes utilizadas.

\section{Grupos de elementos}

\begin{longtable}{p{0.24\linewidth}p{0.34\linewidth}p{0.30\linewidth}}
\toprule
Grupo & Função & Exemplos \\
\midrule
\endfirsthead
\toprule
Grupo & Função & Exemplos \\
\midrule
\endhead
Pré-textuais & Identificar o trabalho e orientar a consulta inicial & Capa, folha de rosto, resumo, listas e sumário \\
Textuais & Desenvolver a investigação e apresentar a argumentação acadêmica & Parte introdutória, desenvolvimento e parte conclusiva \\
Pós-textuais & Registrar fontes e complementar o conteúdo & Referências, glossário, apêndices e anexos \\
\bottomrule
\end{longtable}

\section{Elementos obrigatórios e opcionais}

Alguns elementos são obrigatórios na versão final. Outros são opcionais e devem ser usados quando contribuírem para a identificação, a consulta ou a compreensão do trabalho.

\requisitonormativo

\begin{longtable}{p{0.38\linewidth}p{0.22\linewidth}p{0.28\linewidth}}
\toprule
Elemento & Situação geral & Observação \\
\midrule
\endfirsthead
\toprule
Elemento & Situação geral & Observação \\
\midrule
\endhead
Capa & Obrigatório & Identifica instituição, autoria, título, local e ano. \\
Folha de rosto & Obrigatório & Traz natureza do trabalho, curso e orientação. \\
Dados internacionais de catalogação & Condicionado ao fluxo institucional & A norma prevê sua posição, mas a emissão e a exigência devem seguir a orientação da Biblioteca Central da UnDF. \\
Errata & Opcional & Usada apenas quando há correções posteriores. \\
Folha de aprovação & Obrigatório & Registra aprovação e composição avaliadora. \\
Dedicatória & Opcional & Texto breve e pessoal. \\
Agradecimentos & Opcional & Reconhece contribuições relevantes. \\
Epígrafe & Opcional & Citação breve relacionada ao trabalho. \\
Resumo em português & Obrigatório & Síntese objetiva do trabalho. \\
Resumo em língua estrangeira & Obrigatório & Usualmente em inglês, mas o idioma pode variar por decisão do curso. \\
Listas & Opcional & Usadas quando há quantidade relevante de itens. \\
Sumário & Obrigatório & Último elemento pré-textual. \\
Introdução & Obrigatório & Apresenta tema, problema, objetivos e percurso. \\
Desenvolvimento & Obrigatório & Contém revisão, método, análise e discussão. \\
Parte conclusiva & Obrigatório & Pode receber o título ``Conclusão'' ou ``Considerações finais'', conforme a área e a decisão do curso. \\
Referências & Obrigatório & Lista as fontes citadas. \\
Glossário & Opcional & Define termos técnicos. \\
Apêndice & Opcional & Material elaborado pelo autor. \\
Anexo & Opcional & Material não elaborado pelo autor. \\
Índice & Opcional & Organiza termos de forma sistemática e segue a ABNT NBR 6034:2004. \\
\bottomrule
\end{longtable}

\section{Ordem dos elementos}

A ordem dos elementos deve ser estável. Alterar a posição de resumo, listas, sumário, referências, apêndices ou anexos compromete a padronização e pode dificultar a avaliação.

O sumário deve aparecer ao final dos elementos pré-textuais. A parte textual começa depois do sumário. Referências e demais pós-textuais aparecem após as considerações finais.

A ABNT NBR 14724:2024 define a macroestrutura textual, mas deixa a nomenclatura dos títulos a critério do autor. O trabalho deve ser organizado em seções. Assim, referencial teórico, metodologia, dados, resultados e discussão são formas possíveis de desenvolver o conteúdo, e não títulos universais impostos pela norma.

\section{Coerência estrutural}

A estrutura deve refletir o percurso do trabalho. Se o TCC é pesquisa empírica, a metodologia precisa explicar coleta e análise. Se é pesquisa bibliográfica, a metodologia precisa explicar critérios de busca e seleção. Se é desenvolvimento de produto, o texto deve apresentar requisitos, processo de construção, avaliação e limitações.

\caixaatencao{Não existe estrutura única para todos os cursos e objetos, mas existe uma exigência comum: o leitor precisa entender o caminho seguido e a relação entre problema, método, resultados e conclusão.}
